Следующая теорема является следствием теоремы об образе центра.
\begin{theorem}
  Пусть у облака $[Z]$ нетривиальная стационарная группа, и $Z$
  является его центром. Также, пусть в этом облаке есть пространства
  $Y_{1}, Y_{2}$
  такие, что $\max\big\{ |Y_{1},Z|, |Y_{2}, Z| \big\} = r>0$, а
  $|Y_{1}, Y_{2}|>r$. Тогда, расстояние между облаками $[\Delta_1]$ и $[Z]$
  равно бесконечности.
  \label{thrmDist}
\end{theorem}
\begin{proof}
  У облаков $[\Delta_1]$ и $[Z]$ стационарные группы имеют
  нетривиальное пересечение, и, по лемме \ref{lemmaDist},
  расстояние между ними может быть
  равно либо $0$, либо $\infty$.

  Для доказательства утверждения теоремы
  достаточно будет показать, что расстояние между ними не равно $0$.  Для этого
  необходимо установить, что между ними не может существовать соответствия со
  сколь угодно малым искажением. Итак, пусть $R$ --- соответствие между
  $[\Delta_1]$ и $[Z]$, $\dis R = \varepsilon < \infty$.
  Зафиксируем $Y$ из $R(\Delta_1)$. По теореме \ref{thrmCenterImage}
  расстояние между $Y$ и $Z$ не
  больше $2\varepsilon$.

  По условию теоремы выполнено неравенство:
  $$\max\big\{ |Y_{1},Z|, |Y_{2}, Z| \big\} = r < |Y_{1}, Y_{2}|$$
  Неравенство означает, что существует $c > 0$ такое, что
  $|Y_{1},Y_{2}| = (1 + c)r.$ Вместе с $Y_{1}$ и
  $ Y_{2}$ рассмотрим их прообразы $X_1 \in R^{-1}(Y_{1})$,
  $ X_2 \in R^{-1}(Y_{2})$. Получаем следующую цепочку
  неравенств:
  $$|X_1, \Delta_1| \le |Y_{1}, Y| + \varepsilon \le |Y_{1}, \mathbb{R}|
  + |\mathbb{R}, Y| +\varepsilon \le r + 2\varepsilon + \varepsilon =
  r + 3\varepsilon.$$
  Аналогичное неравенство имеет место для $X_2$, при этом
  $$|X_1, X_2|  \ge |Y_{1}, Y_{2}| - \varepsilon = (1+c)r - \varepsilon.$$
  По замечанию \ref{remUltraMetric}:
  $$|X_1, X_2| \le \max\big\{ |X_1, \Delta_1|, |X_2, \Delta_1| \big\},$$
  $$\Updownarrow$$
  $$(1+c)r - \varepsilon\le r + 3\varepsilon,$$
  $$\Updownarrow$$
  $$\varepsilon \ge \frac{cr}{4}.$$
  Мы получаем оценку снизу для $\varepsilon = \dis R$. Это означает, что
  искажение не может быть произвольно малым, и, следовательно, расстояние между
  пространствами не может быть равно 0. Значит, оно равно бесконечности.
\end{proof}

